\documentclass[12pt,a4paper]{article}

\usepackage[utf8]{inputenc}
\usepackage[T1]{fontenc}
\usepackage[brazilian]{babel}
\usepackage{lmodern}
\usepackage{geometry}
\usepackage{parskip}
\usepackage{array}
\usepackage{booktabs}
\usepackage{longtable}
\usepackage{enumitem}
\usepackage{xcolor}
\usepackage{hyperref}

\geometry{margin=2.5cm}
\setlist[itemize]{leftmargin=*}
\setlength{\LTpre}{0pt}
\setlength{\LTpost}{0pt}
\newcolumntype{L}[1]{>{\raggedright\arraybackslash}p{#1}}
\hypersetup{
  colorlinks=true,
  linkcolor=black,
  urlcolor=blue
}

\begin{document}

\begin{center}
{\Large Relatório para atualização da introdução}\\[0.5em]
{\large Dissertação sobre tempo até o parto, eventos competitivos e estrutura regional em São Paulo}\\[1em]
23 de abril de 2026
\end{center}

\section*{Escopo}

Este relatório foi preparado apenas para a \textbf{introdução} e apenas para a etapa que você definiu agora: atualizar dados e referências e manter o texto coerente com a \textbf{mesma definição operacional} usada no trabalho. Portanto, ele não reescreve ainda a introdução inteira, não mexe no arquivo principal da dissertação e não entra, por enquanto, no ajuste fino de fluidez, estilo ou distribuição final dos parágrafos.

Como a base analítica da dissertação utiliza dados de \textbf{2023}, a lógica adotada neste relatório foi priorizar, sempre que possível, \textbf{fontes recentes publicadas em 2024 ou 2025, mas com dados observados até 2023}. Dados posteriores a 2023 podem até aparecer como contexto complementar, mas não devem ser o eixo principal de comparação com os resultados do estudo.

\section*{Ponto de partida que deve guiar a introdução}

No seu trabalho, o óbito fetal é tratado segundo o critério usado no Brasil para emissão da Declaração de Óbito: \textbf{gestação com duração igual ou superior a 20 semanas, ou peso fetal igual ou superior a 500 g, ou estatura igual ou superior a 25 cm}. Como a sua base não traz estatura, a definição operacional adotada no estudo fica, na prática, centrada em \textbf{idade gestacional} e \textbf{peso fetal}. Esse ponto precisa aparecer de forma simples e clara logo no início da introdução, porque é ele que determina quais indicadores externos podem ou não ser comparados com o seu trabalho.

O mesmo raciocínio vale para os nascidos vivos. O Manual da Declaração de Nascido Vivo define nascimento vivo como a expulsão ou extração completa do produto da concepção que, após a separação, apresente algum sinal de vida, independentemente da duração da gestação. Isso ajuda a amarrar a lógica do trabalho: uma gestação pode terminar em óbito fetal ou em nascimento vivo, e a prematuridade entra justamente dentro do conjunto dos nascidos vivos com idade gestacional menor que 37 semanas.

\section*{Fontes mais úteis para esta etapa}

\small
\begin{longtable}{L{0.20\textwidth} L{0.18\textwidth} L{0.18\textwidth} L{0.16\textwidth} L{0.22\textwidth}}
\toprule
\textbf{Fonte} & \textbf{Tema} & \textbf{Recorte} & \textbf{Aderência ao trabalho} & \textbf{Uso mais indicado na introdução} \\
\midrule
\endhead
Manual da Declaração de Óbito, Ministério da Saúde, 2023 & Definição oficial de óbito fetal no Brasil & Brasil & Alta & Sustentar a definição adotada no estudo e retirar a dependência de definições internacionais logo no começo do texto. \\

Manual da Declaração de Nascido Vivo, Ministério da Saúde, 2022/2023 & Definição de nascido vivo e papel da DNV/Sinasc & Brasil & Alta & Amarrar a presença do nascimento vivo como desfecho concorrente e preparar a entrada da prematuridade na introdução. \\

Rocha et al., 2025, \textit{Revista de Saúde Pública} & Mortalidade fetal com idade gestacional maior ou igual a 20 semanas & Brasil, 1996--2021 & Alta & Principal referência recente para o bloco ``Situação no Brasil''. \\

Plano Nacional de Saúde 2024--2027, atualizado em 2025 & Indicadores oficiais recentes do Ministério da Saúde & Brasil e regiões, com dados até 2021 & Média a alta & Corroborar com fonte oficial do Ministério a persistência das desigualdades regionais, sem depender de dados posteriores à base do estudo. \\

Victor et al., 2025, \textit{BMC Pregnancy and Childbirth} & Tendência nacional da prematuridade & Brasil, 2014--2023 & Alta & Principal referência recente para inserir nascidos vivos e prematuridade na introdução. \\

Boletim Epidemiológico do Ministério da Saúde, 2024 & Perfil epidemiológico dos nascimentos prematuros & Brasil, 2012--2022 & Alta & Fonte oficial complementar para prematuridade, útil se você quiser manter uma referência governamental junto do artigo científico. \\

Painel de Vigilância da Saúde Materna e Perinatal do OOBr, arquivos locais do projeto & Nascidos vivos e prematuridade em anos recentes & Brasil e São Paulo, com prioridade para 2023 & Alta para prematuridade; baixa para mortalidade fetal & Melhor uso para atualizar o parágrafo de São Paulo com números de 2023, mantendo alinhamento temporal com a base do estudo. \\
\bottomrule
\end{longtable}
\normalsize

\section*{Como incorporar na introdução atual}

\subsection*{1. Definições e critérios de classificação do óbito fetal}

\textbf{Trecho atual que precisa ser revisto}

\begin{quote}
Diferentes organizações internacionais adotam critérios distintos de idade gestacional ou de peso para classificar o óbito fetal.
\end{quote}

\textbf{Problema principal}

Hoje essa parte gasta espaço demais com definições internacionais e acaba deixando pouco claro qual definição realmente governa o estudo. Como a sua dissertação não pretende comparar estimativas internacionais sob critérios diferentes, o caminho mais forte é encurtar bastante essa seção e concentrá-la na definição \textbf{brasileira}, que é a mesma usada na construção da base.

\textbf{Como incorporar os dados e referências}

\begin{itemize}
\item Abrir a seção com a definição brasileira do óbito fetal, apoiada no Manual da Declaração de Óbito do Ministério da Saúde.
\item Explicar em seguida, em uma frase curta, que a base do estudo não contém estatura e, por isso, a operacionalização adotada considera idade gestacional e peso fetal.
\item Inserir logo depois a definição de nascido vivo a partir do Manual da DNV, para que o leitor entenda desde cedo por que o nascimento vivo aparece no trabalho como evento concorrente.
\item Reduzir a discussão sobre OMS, CID-10 e definições internacionais a uma observação breve, apenas para dizer que elas existem e que, por isso, comparações externas precisam de cautela.
\end{itemize}

\textbf{Trecho que pode entrar depois da revisão}

Pode entrar algo nesta linha:

\begin{quote}
Neste trabalho, o óbito fetal é considerado segundo o critério adotado no Brasil para emissão da Declaração de Óbito, isto é, perdas gestacionais com 20 semanas ou mais, ou peso fetal de 500 g ou mais, ou estatura de 25 cm ou mais. Como a base utilizada não dispõe de informação sobre estatura, a definição operacional adotada considerou a idade gestacional e o peso fetal. Em contraste, o nascimento vivo corresponde à expulsão ou extração completa do produto da concepção que, após a separação do corpo materno, apresenta qualquer sinal de vida, independentemente da duração da gestação. Essa distinção é central para o presente estudo, porque o acompanhamento gestacional pode terminar em óbito fetal ou em nascimento vivo, sendo este último também relevante para o estudo da prematuridade.
\end{quote}

\textbf{Fontes prioritárias deste bloco}

\begin{itemize}
\item Manual da Declaração de Óbito, Ministério da Saúde, 2023.
\item Manual da Declaração de Nascido Vivo, Ministério da Saúde, 2022/2023.
\end{itemize}

\subsection*{2. Panorama global}

\textbf{Trecho atual que precisa ser revisto}

\begin{quote}
Estimativas recentes indicam que aproximadamente 1,9 milhão de bebês nasceram mortos em 2023, considerando gestações com 28 semanas ou mais.
\end{quote}

\textbf{Problema principal}

Esse bloco usa uma definição internacional diferente da sua. Por isso, ele pode continuar existindo, mas só como \textbf{contexto muito breve}, sem ocupar espaço demais e sem dar a impressão de que os valores globais são diretamente comparáveis aos seus resultados ou aos indicadores brasileiros baseados em 20 semanas.

\textbf{Como incorporar os dados e referências}

\begin{itemize}
\item Manter, no máximo, um parágrafo curto de contexto global.
\item Deixar explícito que os números globais mais recentes costumam usar corte de 28 semanas para fins de comparabilidade internacional.
\item Evitar fazer do bloco global o centro da introdução, já que o restante do trabalho está apoiado na definição brasileira.
\end{itemize}

\textbf{Observação importante}

Se você quiser seguir de forma ainda mais estrita a orientação da sua orientadora, este bloco pode até ser reduzido a duas ou três frases e perder o papel de seção autônoma. Do ponto de vista de coerência conceitual, isso faria sentido.

Além disso, como o foco comparativo do estudo está em dados até 2023, este bloco global deve funcionar apenas como pano de fundo e não como referência principal para confronto dos achados.

\subsection*{3. Situação no Brasil}

\textbf{Trecho atual que precisa ser revisto}

\begin{quote}
Em 2015, a taxa de natimortalidade brasileira foi estimada em aproximadamente 10,8 óbitos fetais por 1000 nascimentos.
\end{quote}

\textbf{Problema principal}

O trecho atual depende demais de números antigos e alterna definições sem deixar isso claro. Aqui você já tem, hoje, uma referência muito melhor: o artigo de Rocha et al. de 2025, que trabalha com \textbf{idade gestacional maior ou igual a 20 semanas}, ou seja, dialoga muito melhor com a base do seu estudo.

\textbf{Dados recentes que podem entrar}

\begin{itemize}
\item Rocha et al. analisaram o Brasil de 1996 a 2021 com corte em 20 semanas e encontraram taxa global de mortalidade fetal de \textbf{11,4 por 1000 nascimentos} nesse período.
\item No mesmo estudo, o país registrou \textbf{29.325 óbitos fetais em 2021} no estrato de 20 semanas ou mais.
\item Cerca de \textbf{93\% dos óbitos} ocorreram antes do parto, o que ajuda a contextualizar a relevância obstétrica do problema.
\item O artigo também mostra persistência de desigualdades regionais, com piores cenários concentrados no Norte, Nordeste e Centro-Oeste.
\item No Plano Nacional de Saúde 2024--2027, atualizado em 2025, o Ministério da Saúde informa, para 2021, taxa de \textbf{10 óbitos fetais por 1000 nascidos vivos} no Brasil, com valores mais altos no Norte e Nordeste (\textbf{11,5}) e mais baixos no Sul (\textbf{7,5}), o que reforça a leitura de desigualdade regional.
\end{itemize}

\textbf{Como incorporar os dados de nascidos vivos e prematuridade neste mesmo bloco}

Como a sua introdução hoje está muito centrada em óbito fetal, vale acoplar ainda neste bloco uma passagem curta sobre nascimento vivo e prematuridade, para preparar melhor o leitor para o desenho do estudo.

\begin{itemize}
\item Victor et al. analisaram mais de \textbf{25,5 milhões de nascidos vivos} no Brasil entre 2014 e 2023.
\item Nesse estudo, a prevalência nacional de prematuridade subiu de \textbf{11,3\% em 2014} para \textbf{11,9\% em 2023}.
\item As maiores frequências foram observadas em mães com 35 anos ou mais, em mães com menos de 20 anos e nas regiões Norte e Nordeste.
\end{itemize}

\textbf{Trecho que pode entrar depois da revisão}

Pode entrar algo nesta linha:

\begin{quote}
No Brasil, o óbito fetal continua sendo um problema relevante de saúde pública mesmo após a redução observada nas últimas décadas. Em estudo nacional recente com dados do SIM e definição de óbito fetal a partir de 20 semanas de gestação, Rocha et al. observaram taxa global de 11,4 óbitos fetais por 1000 nascimentos entre 1996 e 2021, com 29.325 registros em 2021 e predominância de ocorrências antes do parto. Além da magnitude do problema, persistem desigualdades regionais, com maiores taxas e reduções mais lentas no Norte, Nordeste e Centro-Oeste. Em paralelo, os nascidos vivos também exigem atenção, especialmente pela frequência da prematuridade. Em análise nacional do Sinasc, a prevalência de parto prematuro aumentou de 11,3\% em 2014 para 11,9\% em 2023, mantendo-se mais elevada nos extremos de idade materna e em regiões socialmente mais vulneráveis.
\end{quote}

\textbf{Fontes prioritárias deste bloco}

\begin{itemize}
\item Rocha JBF, Bezerra IMP, Oliveira EKS, Sena ABE, Leitão FNC, Abreu LC. \textit{Space-time trends in fetal mortality in Brazil, 1996--2021}. Revista de Saúde Pública, 2025.
\item Ministério da Saúde. \textit{Plano Nacional de Saúde 2024--2027}. Atualizado em maio de 2025.
\item Victor A, Nhancololo AM, Bah HAF et al. \textit{National and regional temporal trends and forecasting of preterm birth in Brazil: evidence from national birth data (2014--2023) with projections to 2030}. BMC Pregnancy and Childbirth, 2025.
\item Ministério da Saúde. \textit{Boletim Epidemiológico, volume 55, número 13: Perfil epidemiológico dos nascimentos prematuros no Brasil, 2012 a 2022}. 2024.
\end{itemize}

\subsection*{4. Contexto no estado de São Paulo}

\textbf{Trecho atual que precisa ser revisto}

\begin{quote}
Entre 2010 e 2014, por exemplo, a taxa de mortalidade fetal no estado de São Paulo foi de aproximadamente 7,9 óbitos fetais por 1000 gestações.
\end{quote}

\textbf{Problema principal}

O problema aqui não é apenas a idade da referência. É também o fato de que o parágrafo paulista ficou preso ao óbito fetal e não conversa com o restante do estudo, que também olha para nascimento vivo e prematuridade. Para este bloco, os dados mais fortes que você já pode usar agora são os do \textbf{Painel de Vigilância da Saúde Materna e Perinatal} disponível na pasta de suporte do projeto, ao menos para \textbf{nascidos vivos} e \textbf{prematuridade}.

\textbf{Dados recalculados a partir do painel local}

Recalculando os arquivos do painel presentes no projeto, obtêm-se os seguintes valores:

\begin{itemize}
\item \textbf{Brasil, 2023}: 2.537.511 nascidos vivos e prematuridade de 11,96\%.
\item \textbf{São Paulo, 2023}: 503.903 nascidos vivos e prematuridade de 12,02\%.
\end{itemize}

Para este relatório, os valores de 2024 foram deixados de fora da proposta principal de incorporação, porque o objetivo aqui é manter o texto da introdução comparável com uma pesquisa cuja base termina em 2023.

\textbf{Como incorporar esses dados}

\begin{itemize}
\item Substituir o parágrafo centrado em 2010--2014 por um parágrafo que apresente São Paulo como um cenário de alto volume de nascidos vivos e, ao mesmo tempo, com carga relevante de prematuridade.
\item Usar esses números para aproximar a introdução do problema real do seu estudo, que envolve tanto óbito fetal quanto nascimento vivo e parto prematuro.
\item Evitar, por enquanto, usar no texto da introdução a taxa de mortalidade fetal do painel do OOBr como se fosse diretamente equivalente ao seu trabalho.
\item Evitar fazer do ano de 2024 o eixo do argumento descritivo, já que a comparação principal da dissertação deve permanecer ancorada até 2023.
\end{itemize}

\textbf{Por que ter cautela com a mortalidade fetal do painel}

Na documentação interna do painel, a taxa de mortalidade fetal aparece baseada em \textbf{óbitos fetais com 22 semanas ou mais}, e a fórmula documentada não coincide perfeitamente com a operacionalização do seu estudo. Então, para o texto da dissertação, o uso mais seguro neste momento é:

\begin{itemize}
\item usar o painel para \textbf{nascidos vivos} e \textbf{prematuridade}; e
\item deixar a atualização da mortalidade fetal paulista para quando você optar entre uma fonte pública recente com definição explícita compatível ou o cálculo direto a partir da sua própria base SIM-DOFET/SINASC.
\end{itemize}

\textbf{Trecho que pode entrar depois da revisão}

Pode entrar algo nesta linha:

\begin{quote}
No estado de São Paulo, a discussão também precisa considerar o grande volume de nascidos vivos e a persistência da prematuridade. A partir dos dados do Painel de Vigilância da Saúde Materna e Perinatal disponíveis no projeto, o estado registrou 503.903 nascidos vivos em 2023, com prematuridade de 12,02\% entre os nascidos vivos. Esses valores mostram que, mesmo em um estado com ampla rede assistencial, a distribuição dos desfechos gestacionais continua exigindo análise conjunta do nascimento vivo, da prematuridade e do óbito fetal, sobretudo quando se considera a heterogeneidade regional interna.
\end{quote}

\textbf{Fontes prioritárias deste bloco}

\begin{itemize}
\item Painel de Vigilância da Saúde Materna e Perinatal do Observatório Obstétrico Brasileiro, com recálculo a partir dos arquivos CSV presentes na pasta de suporte do projeto.
\item Victor et al., 2025, como base científica complementar para o enquadramento da prematuridade.
\end{itemize}

\section*{Seções da introdução que não precisam ser mexidas agora}

Pelos critérios desta etapa, eu não mexeria ainda de forma pesada em:

\begin{itemize}
\item \textbf{Fatores associados}, porque a maior necessidade imediata ali não é atualizar definição, e sim ajustar tom, fluidez e seleção bibliográfica.
\item \textbf{Abordagens analíticas e relevância metodológica}, porque essa parte conversa mais com o desenho estatístico do estudo do que com a atualização de indicadores descritivos.
\end{itemize}

\section*{Síntese prática para a próxima versão da introdução}

Se a atualização for feita agora apenas com base neste relatório, a introdução já deve ganhar quatro melhorias concretas:

\begin{itemize}
\item a definição usada no trabalho ficará clara desde o início;
\item os dados brasileiros deixarão de depender de referências antigas e de critérios misturados;
\item a introdução passará a incorporar nascidos vivos e prematuridade de forma coerente com o objetivo do estudo;
\item o trecho sobre São Paulo deixará de depender de números muito antigos e passará a dialogar melhor com a base, sem deslocar a comparação para anos posteriores a 2023.
\end{itemize}

\section*{Referências citadas neste relatório}

\begin{enumerate}[leftmargin=*]
\item Brasil. Ministério da Saúde. \textit{Declaração de Óbito: manual de instruções para preenchimento}. Brasília: Ministério da Saúde; 2023. Disponível em: \href{https://www.gov.br/saude/pt-br/centrais-de-conteudo/publicacoes/svsa/vigilancia/declaracao-de-obito-manual-de-instrucoes-para-preenchimento.pdf/view}{gov.br/saude/.../declaracao-de-obito-manual-de-instrucoes-para-preenchimento}.

\item Brasil. Ministério da Saúde. \textit{Declaração de Nascido Vivo: manual de instruções para preenchimento}. 4.\ ed. Brasília: Ministério da Saúde; 2022. Disponível em: \href{https://bvsms.saude.gov.br/bvs/publicacoes/declaracao_nascido_vivo_manual_4ed.pdf}{bvsms.saude.gov.br/.../declaracao\_nascido\_vivo\_manual\_4ed.pdf}.

\item Rocha JBF, Bezerra IMP, Oliveira EKS, Sena ABE, Leitão FNC, Abreu LC. Space-time trends in fetal mortality in Brazil, 1996--2021. \textit{Revista de Saúde Pública}. 2025;59:e2. DOI: \url{https://doi.org/10.11606/s1518-8787.2025059006194}.

\item Brasil. Ministério da Saúde. \textit{Plano Nacional de Saúde 2024--2027}. Atualizado em maio de 2025. Disponível em: \href{https://bvsms.saude.gov.br/bvs/publicacoes/plano_nacional_saude_pns_2024_2027.pdf}{bvsms.saude.gov.br/.../plano\_nacional\_saude\_pns\_2024\_2027.pdf}.

\item Victor A, Nhancololo AM, Bah HAF, Xavier SP, Gotine AREM, Moreira VA, et al. National and regional temporal trends and forecasting of preterm birth in Brazil: evidence from national birth data (2014--2023) with projections to 2030. \textit{BMC Pregnancy and Childbirth}. 2025;25:1339. DOI: \url{https://doi.org/10.1186/s12884-025-08536-6}.

\item Brasil. Ministério da Saúde. \textit{Boletim Epidemiológico. Volume 55, número 13: Perfil epidemiológico dos nascimentos prematuros no Brasil, 2012 a 2022}. 2024. Disponível em: \href{https://www.gov.br/saude/pt-br/centrais-de-conteudo/publicacoes/boletins/epidemiologicos/edicoes/2024/boletim-epidemiologico-volume-55-no-13.pdf/view}{gov.br/saude/.../boletim-epidemiologico-volume-55-no-13}.

\item Observatório Obstétrico Brasileiro. \textit{Painel de Vigilância da Saúde Materna e Perinatal}. Consulta apoiada pelos arquivos do painel presentes na pasta de suporte do projeto, com recálculo próprio dos indicadores de nascidos vivos e prematuridade para Brasil e São Paulo em 2023. Página institucional: \href{https://observatorioobstetricobr.org/paineis/}{observatorioobstetricobr.org/paineis}.
\end{enumerate}

\end{document}
